\documentclass[10pt]{exam}
\usepackage[phy]{template-for-exam}
\usepackage{enumitem,multicol}
\setlist[itemize]{itemsep=-1ex,partopsep=1ex,parsep=1ex}

\title{Water Rockets}
\author{Rohrbach}
\date{\today}

\begin{document}
\maketitle


\section*{Launches}

Describe the motion (acceleration, height, time in the air) of each rocket.

\paragraph{Launch \#1} Pressure: 60~PSI; Mass: about 1~kg
\vs 

\paragraph{Launch \#2} Pressure: 90~PSI; Mass: about 1~kg
\vs 

\paragraph{Launch \#3} Pressure: 60~PSI; Mass: about 1.7~kg
\vs 

\section*{Questions}

\begin{questions}
  \question
    Use Newton's Third Law to explain how the rocket works.
    \vs

  \question
    Notice that the thrust from the rocket ends pretty quickly.  However, the rocket does not immediately stop its upward motion.  Explain why this is using Newton's Laws.
    \vs

  \question
    Compare the acceleration of Launch \#1 to Launch \#2.  Explain which one had more acceleration using Newton's Laws.
    \vs

  \question
    What happened in Launch \#3?  Why do you think its acceleration was so much different than the other two?
    \vs

  \pagebreak

  \question
    Explain why a rocket should have a cone on the top of it.  Relate your explanation to at least one of Newton's Laws.
    \vs

  \question
    Above, you used Newton's Third Law to explain how the rocket launches.  Considering that the upward force on the rocket is equal and opposite to the downward force on the water, why do these two forces not cancel each other out?
    \vs
  
\uplevel{\section*{Practice Calculations}}


  \question
    A 1.2-kg water rocket experiences 130~N of upward thrust and 95~N of air drag.  What is the rocket's acceleration?  (\emph{Hint:} Your free-body diagram should have three forces.)
    \vs[3]

  \question
    Once the rocket reaches its maximum height, its mass is only 0.13~kg.  As it falls, it accelerates downward at 6.5~m/s$^2$.  Calculate the air drag acting on the rocket.  (\emph{Hint:} Your free-body diagram should have two forces.  Also, think carefully about which way the air drag force points.)
    \vs[3]

\end{questions}

\end{document}